\begin{helper}
Let \(U\) be a finite terminal coordinate set.  Let \(\mathcal P\) be a
finite set of transcript labels, and let
\(\operatorname{Realized}(\beta,\tau)\) be a predicate on branch labels and
transcripts.  For every realized pair with \(\tau\in\mathcal P\), let
\(\mathcal R_{\beta,\tau}\) be a finite set of released terminal
assignments, and let \(\rho_{\beta,\tau}\) be its released cylinder mass.
Let \(m(A)\) be a nonnegative assignment mass on complete terminal
assignments \(A\subseteq U\), with
\[
\sum_{A\subseteq U}m(A)\le 1 .
\]
Assume:
\[
\mathcal R_{\beta,\tau}\subseteq \{A:A\subseteq U\}
\]
for every realized \((\beta,\tau)\) with \(\tau\in\mathcal P\);
\[
\rho_{\beta,\tau}=\sum_{A\in\mathcal R_{\beta,\tau}}m(A)
\]
for every such realized pair; and distinct realized pairs have disjoint
released assignment sets, in the following pointwise form: if
\((\beta,\tau)\ne(\beta',\tau')\), both pairs are realized, both transcripts
belong to \(\mathcal P\), and \(A\in\mathcal R_{\beta,\tau}\), then
\[
A\notin\mathcal R_{\beta',\tau'} .
\]
Then the released cylinders form a subprobability:
\[
\sum_{\beta}
  \sum_{\tau\in\mathcal P:\operatorname{Realized}(\beta,\tau)}
    \rho_{\beta,\tau}
\le 1 .
\]
\end{helper}
\begin{proof}
Let
\[
\mathcal C
  =\{(\beta,\tau):\tau\in\mathcal P,\,
      \operatorname{Realized}(\beta,\tau)\}.
\]
The displayed double sum over \(\beta\) and realized transcripts is exactly
the finite sum of \(\rho_{\beta,\tau}\) over \(\mathcal C\).  By the exact
released-mass hypothesis, this is
\[
\sum_{(\beta,\tau)\in\mathcal C}
  \sum_{A\in\mathcal R_{\beta,\tau}}m(A).
\]

The pointwise disjointness hypothesis says that the family
\(\{\mathcal R_{\beta,\tau}\}_{(\beta,\tau)\in\mathcal C}\) is pairwise
disjoint.  Therefore finite additivity over a disjoint finite union gives
\[
\sum_{(\beta,\tau)\in\mathcal C}
  \sum_{A\in\mathcal R_{\beta,\tau}}m(A)
=
\sum_{A\in\bigcup_{(\beta,\tau)\in\mathcal C}
  \mathcal R_{\beta,\tau}}m(A).
\]
Every released assignment lies in the complete assignment space
\(\{A:A\subseteq U\}\).  Since all masses \(m(A)\) on complete terminal
assignments are nonnegative,
\[
\sum_{A\in\bigcup_{(\beta,\tau)\in\mathcal C}
  \mathcal R_{\beta,\tau}}m(A)
\le
\sum_{A\subseteq U}m(A)
\le 1 .
\]
Combining these equalities and inequalities proves the claimed
subprobability bound.
\end{proof}
